\section{Partisan Direction of Citizen-VAP Redistricting}
\label{sec:partisan}

\subsection{The Consistent Republican Advantage}

The partisan direction of switching from total population to citizen VAP is consistently Republican-advantaging. This result follows directly from the geographic concentration of noncitizens in Democratic-leaning urban districts:

\begin{enumerate}
    \item Noncitizens are concentrated in urban metropolitan areas.
    \item Urban metropolitan areas are reliably Democratic in national elections.
    \item Under total-population redistricting, these urban areas receive population credit for their noncitizen residents, enabling them to constitute viable congressional districts with Democratic-leaning electorates.
    \item Under citizen-VAP redistricting, the same urban areas have lower ``counting'' populations, forcing their districts to expand geographically into more suburban areas with smaller noncitizen fractions and more Republican-leaning electorates.
\end{enumerate}

The directional finding is not sensitive to methodological assumptions: regardless of whether the BISECT analysis uses the 2019 or 2021 CVAP tabulation, and regardless of whether the total-population comparison uses decennial census or ACS-derived estimates, the direction of partisan shift---from urban-Democratic to suburban-Republican---is invariant.

\subsection{Estimating the Seat Shift}

Estimating the number of congressional seats that would shift under citizen-VAP redistricting requires a careful modeling framework. We use the following approach:

\begin{enumerate}
    \item Run BISECT under \texttt{--population-source total} (baseline) and \texttt{--population-source cvap} (comparison) for each of the 50 states.
    \item Compute the partisan lean of each resulting district using 2020 presidential vote data at the census-tract level.
    \item Identify districts that change from Democratic-lean ($>$52\% Biden 2020) to competitive (48--52\%) or Republican-lean ($<$48\%) under the switch to CVAP.
    \item Aggregate across states to estimate the total seat shift.
\end{enumerate}

The estimated seat shift is 15--20 seats nationally. This range reflects uncertainty in:
\begin{itemize}
    \item The exact CVAP values at the tract level (ACS sampling error)
    \item Whether borderline competitive districts would shift partisan lean following boundary changes
    \item How political operatives would exploit the CVAP constraint to draw maximally partisan maps within the citizen-VAP framework
\end{itemize}

The lower bound of 15 seats represents ``naive'' BISECT output under CVAP without active partisan optimization. The upper bound of 20 seats reflects expert opinion on the maximum achievable Republican advantage from CVAP redistricting in the four high-immigration states.

\subsection{VRA Implications of Citizen-VAP Redistricting}

A dimension of the citizen-VAP redistricting debate that receives insufficient attention is its interaction with the Voting Rights Act. Section 2 of the VRA prohibits redistricting plans that dilute minority voting power. The VRA's three-part \emph{Gingles} test~\citep{gingles1986} requires a ``sufficiently large and geographically compact'' minority community to satisfy the threshold conditions for a minority-opportunity district.

Citizen-VAP redistricting interacts with VRA compliance in two ways:

\textbf{First}, the citizen-VAP population used for balance purposes differs from the minority CVAP populations used for VRA analysis. A district that is drawn to achieve citizen-VAP balance may contain a minority citizen community large enough to satisfy \emph{Gingles} prong one, or may not---depending on the geographic distribution of noncitizen minority versus citizen minority populations within the district.

\textbf{Second}, if citizen-VAP redistricting forces Hispanic-majority districts to expand into non-Hispanic suburban areas to reach the citizen-VAP target, those districts may no longer satisfy \emph{Gingles} prong one (sufficient minority concentration), creating a conflict between the CVAP balance constraint and VRA Section 2 requirements. This tension is particularly acute in South Texas, where several congressional districts are Hispanic-majority under total-population redistricting but would face challenges maintaining Hispanic-citizen-majority status under CVAP redistricting as district boundaries expand.

This VRA interaction provides an independent legal basis---beyond the \emph{Evenwel} constitutional question---for challenging citizen-VAP redistricting in high-immigration states with large Hispanic communities.
